\chapter{Sprint 1 : Authentification et gestion des profils}
\minitoc
\newpage

\section{Sprint I : Authentification et gestion des profils}

\subsection{Objectif du Sprint}
Le premier Sprint constitue la pierre angulaire de toute la plateforme GOSPORT. Sans une infrastructure technique solide et un système d'authentification robuste, aucune des fonctionnalités métier ne peut être construite de manière fiable. Ce sprint répond à la question fondamentale : \textit{Comment permettre à un utilisateur de rejoindre la plateforme de manière sécurisée et de configurer son espace personnalisé ?}

Ce premier incrément livre trois réalisations critiques : l'initialisation complète de la base de données relationnelle PostgreSQL avec TypeORM, la mise en place de la sécurité applicative (hachage, jetons JWT, OAuth 2.0), et la création des formulaires d'inscription et d'onboarding sur l'application web Angular.

\subsection{Backlog du Sprint I}

\begin{table}[H]
\centering
\renewcommand{\arraystretch}{1.4}
\resizebox{\textwidth}{!}{%
\begin{tabular}{|c|c|p{5cm}|p{6.5cm}|c|}
\hline
\rowcolor[HTML]{DAE8FC}
\textbf{ID} & \textbf{Acteur} & \textbf{User Story} & \textbf{Tâche Technique} & \textbf{Priorité} \\ \hline
1.1 & Tous & Créer un compte sécurisé & Entité User, hachage bcrypt, route POST /register & Must \\ \hline
1.2 & Tous & S'authentifier avec email/mdp & Route POST /login, génération token JWT & Must \\ \hline
1.3 & Tous & Se connecter via Google OAuth & Intégration Passport.js OAuth 2.0 & Must \\ \hline
1.4 & Athlète & Compléter le profil biométrique & Formulaire onboarding Angular, entité Athlete & Must \\ \hline
1.5 & Coach & Configurer son profil professionnel & Entité Coach, formulaire profil coach & Should \\ \hline
\end{tabular}%
}
\caption{Backlog du Sprint 1 -- Authentification et Onboarding}
\end{table}

\section{Analyse}

\subsection{Diagramme de cas d'utilisation du premier sprint}
\begin{figure}[H]
\centering
\includegraphics[width=0.88\textwidth]{Sprint1_Cas_Utilisation.png}
\caption{Diagramme de cas d'utilisation -- Sprint 1}
\end{figure}

\subsection{Tableau descriptif du cas d'utilisation « S'inscrire »}

\subsection{Tableau descriptif du cas d'utilisation « Se connecter »}

\subsection{Tableau descriptif du cas d'utilisation « Récupérer le mot de passe »}

\subsection{Tableau descriptif du cas d'utilisation « Compléter le profil Inspecteur »}

\subsection{Tableau descriptif du cas d'utilisation « Gérer le profil Inspecteur »}

\subsection{Tableau descriptif du cas d'utilisation « Compléter le profil Enseignant »}

\subsection{Tableau descriptif du cas d'utilisation « Gérer le profil Enseignant »}

\section{Conception}

\subsection{Diagrammes de séquences}

\subsection{A) Diagramme de séquence « Processus d'inscription utilisateur »}
\begin{figure}[H]
\centering
\includegraphics[width=0.88\textwidth]{Sprint1_Sequence_Auth.png}
\caption{Diagramme de séquence -- Processus d'inscription utilisateur}
\end{figure}

\subsection{B) Diagramme de séquence « Processus de connexion utilisateur »}

\subsection{C) Diagramme de séquence « Configuration du profil Enseignant »}

\subsection{D) Diagramme de séquence « Configuration du profil Inspecteur »}

\subsection{Diagramme de classes du premier sprint}
\begin{figure}[H]
\centering
\includegraphics[width=0.9\textwidth]{Sprint1_Classes.png}
\caption{Diagramme de classes -- Authentification et Profils}
\end{figure}

Les entités clés de ce sprint sont : \texttt{User} (rôle énuméré : ATHLETE, COACH, ADMIN), \texttt{Athlete} (profil biométrique : poids, taille, objectif) et \texttt{Coach} (spécialité, tarif, disponibilité). La relation \texttt{OneToOne} entre \texttt{User} et ses profils spécialisés permet une gestion claire des droits d'accès et des données métier.

\section{Réalisation}

\subsection{Description des interfaces}

\section*{Conclusion du Sprint I}
\addcontentsline{toc}{section}{Conclusion du Sprint I}
\mtcaddsection
\markboth{Conclusion du Sprint I}{}
À l'issue de ce premier sprint, les fondations techniques et sécuritaires de GOSPORT sont opérationnelles. L'authentification locale et sociale, la persistance des données via TypeORM, la configuration des Guards Angular et le parcours d'onboarding sont livrés et validés. Ces fondations permettent aux sprints suivants de se concentrer sur les fonctionnalités métier.
